% META: {"id": "TEXP-ARI-CO-026", "chapitre": "TEXP-ARITHMETIQUE", "type_objet": "corrige", "exercice_ref": "TEXP-ARI-EX-026", "capacites_codes": ["C8"], "sources_inspiration": [], "status": "generated"}
% BEGIN-VERIFY
% assert pow(2, 340, 341) == 1
% assert 341 == 11 * 31
% assert pow(3, 340, 341) == 56
% assert pow(3, 340, 341) != 1
% END-VERIFY
\begin{corrige}{TEXP-ARI-EX-026}
\textbf{1.} Un calcul d'exponentiation modulaire donne $2^{340}\equiv1\ [341]$.

\textbf{2.} $341=11\times31$ : il n'est pas premier.

\textbf{3.} La réciproque est fausse : passer le test de Fermat en base $2$ ne
prouve pas la primalité. Un tel nombre est dit pseudo-premier de base $2$. En
base $3$, en revanche, $3^{340}\equiv56\not\equiv1\ [341]$ : la base $3$ est un
témoin de non-primalité.
\end{corrige}
